% GENERATED FILE. Edit canonical lessons, then run npm run generate:books.
% canonical-source-hash: fnv1a64:bf153c34accc2a0c
% canonical-lessons: MR-C295-nuksan, MR-C295-durghatna, MR-C295-apatti, MR-C295-sankat, MR-C295-bhisan, MR-R295-review-288-295

\chapter{After the Disaster}
\label{ch:mr-after-the-disaster}
\hlchaptermodality{295}

\begin{chapteropening}
\textbf{By the end of this chapter:} \emph{I can talk about the damage a disaster leaves behind.}

Complete the last lesson of chapter 295: I can talk about the damage a disaster leaves behind.
\end{chapteropening}

\section[\texorpdfstring{\mr{नुकसान} (nuksān)}{नुकसान (nuksān)}]{\mr{नुकसान} (nuksān) — damage; a loss}

\label{lesson:MR-C295-nuksan}



\begin{quote}
Before the new one: say the Marathi for a volunteer, then the Marathi for to get stuck.
\end{quote}

\subsection*{You'll want to know: \mr{नुकसान}}
\textbf{\mr{नुकसान}} — \emph{nuksān} — \textquotedblleft{}damage; a loss\textquotedblright{}.

\begin{quote}
\mr{जर} \mr{नुकसान} \mr{मोठं} \mr{असेल}, \mr{तर} \mr{गावाला} \mr{मदत} \mr{मिळेल}.
\end{quote}

\begin{quote}
\emph{jar nuksān moṭhaṁ asel, tar gāvālā madat miḷel.}
\end{quote}

\begin{quote}
If the damage is great, the village will get help.
\end{quote}

Neuter.

\mr{नुकसान} \mr{होणे}, a loss is suffered; \mr{जर} … \mr{असेल} \mr{तर}, if it is, for a state.

\begin{cousinweb}
\textbf{\mr{नुकसान}} (\emph{nuksān}) is a loan from Persian \emph{nuqsān}, usually traced to Arabic \emph{nuqṣān}, 'a shortfall, a deficit'.
\end{cousinweb}

\begin{grammarlens}[title={What you've built}]
One more word for telling what happened.
\end{grammarlens}

\subsection*{Your turn}
\begin{itemize}
\raggedright
  \item \emph{Say it:} \emph{nuksān}
  \item \emph{Say it:} \emph{nuksān}, once more
  \item \emph{Say it:} \emph{jar nuksān moṭhaṁ asel, tar gāvālā madat miḷel.} — If the damage is great, the village will get help.
  \item \emph{Recall:} say the Marathi for a volunteer, then the Marathi for to get stuck, then say \emph{nuksān} again
\end{itemize}

\begin{tcolorbox}[breakable,colback=teal!4,colframe=teal!35!black,title={Before you move on}]
What does \mr{नुकसान} mean? (\textquotedblleft{}Damage; a loss\textquotedblright{}.) Say it once more.
\end{tcolorbox}
\section[\texorpdfstring{\mr{दुर्घटना} (durghaṭnā)}{दुर्घटना (durghaṭnā)}]{\mr{दुर्घटना} (durghaṭnā) — a disaster; a terrible accident}

\label{lesson:MR-C295-durghatna}



\begin{quote}
Before the new one: say the Marathi for to get stuck, then the Marathi for damage.
\end{quote}

\subsection*{You'll want to know: \mr{दुर्घटना}}
\textbf{\mr{दुर्घटना}} — \emph{durghaṭnā} — \textquotedblleft{}a disaster; a terrible accident\textquotedblright{}.

\begin{quote}
\mr{त्या} \mr{दुर्घटनेत} \mr{अनेक} \mr{घरं} \mr{वाहून} \mr{गेली}.
\end{quote}

\begin{quote}
\emph{tyā durghaṭnet anek gharaṁ vāhūn gelī.}
\end{quote}

\begin{quote}
Many houses were swept away in that disaster.
\end{quote}

Feminine.

Built from \mr{दुर्}-, ill, and \mr{घटना}, an event.

\begin{cousinweb}
\textbf{\mr{दुर्घटना}} (\emph{durghaṭnā}) is a learned compound built from the Sanskrit prefix \emph{dur-}, 'bad, ill', and \emph{ghaṭanā}, 'a happening' (Monier-Williams).
\end{cousinweb}

\begin{grammarlens}[title={What you've built}]
One more word for telling what happened.
\end{grammarlens}

\subsection*{Your turn}
\begin{itemize}
\raggedright
  \item \emph{Say it:} \emph{durghaṭnā}
  \item \emph{Say it:} \emph{durghaṭnā}, once more
  \item \emph{Say it:} \emph{tyā durghaṭnet anek gharaṁ vāhūn gelī.} — Many houses were swept away in that disaster.
  \item \emph{Recall:} say the Marathi for to get stuck, then the Marathi for damage, then say \emph{durghaṭnā} again
\end{itemize}

\begin{tcolorbox}[breakable,colback=teal!4,colframe=teal!35!black,title={Before you move on}]
What does \mr{दुर्घटना} mean? (\textquotedblleft{}A disaster; a terrible accident\textquotedblright{}.) Say it once more.
\end{tcolorbox}
\section[\texorpdfstring{\mr{आपत्ती} (āpattī)}{आपत्ती (āpattī)}]{\mr{आपत्ती} (āpattī) — a calamity}

\label{lesson:MR-C295-apatti}



\begin{quote}
Before the new one: say the Marathi for damage, then the Marathi for a disaster.
\end{quote}

\subsection*{You'll want to know: \mr{आपत्ती}}
\textbf{\mr{आपत्ती}} — \emph{āpattī} — \textquotedblleft{}a calamity\textquotedblright{}.

\begin{quote}
\mr{अशा} \mr{आपत्तीत} \mr{सगळे} \mr{एकमेकांना} \mr{मदत} \mr{करतात}.
\end{quote}

\begin{quote}
\emph{aśā āpattīt sagḷe ekmekānnā madat kartāt.}
\end{quote}

\begin{quote}
In such a calamity everyone helps one another.
\end{quote}

Feminine.

\mr{नैसर्गिक} \mr{आपत्ती}, a natural calamity.

\begin{cousinweb}
\textbf{\mr{आपत्ती}} (\emph{āpattī}) is a tatsama, taken from Sanskrit \emph{āpatti}, 'a happening, an occurrence' (Monier-Williams); Marathi keeps it for the happenings that bring ruin.
\end{cousinweb}

\begin{grammarlens}[title={What you've built}]
One more word for telling what happened.
\end{grammarlens}

\subsection*{Your turn}
\begin{itemize}
\raggedright
  \item \emph{Say it:} \emph{āpattī}
  \item \emph{Say it:} \emph{āpattī}, once more
  \item \emph{Say it:} \emph{aśā āpattīt sagḷe ekmekānnā madat kartāt.} — In such a calamity everyone helps one another.
  \item \emph{Recall:} say the Marathi for damage, then the Marathi for a disaster, then say \emph{āpattī} again
\end{itemize}

\begin{tcolorbox}[breakable,colback=teal!4,colframe=teal!35!black,title={Before you move on}]
What does \mr{आपत्ती} mean? (\textquotedblleft{}A calamity\textquotedblright{}.) Say it once more.
\end{tcolorbox}
\section[\texorpdfstring{\mr{संकट} (saṅkaṭ)}{संकट (saṅkaṭ)}]{\mr{संकट} (saṅkaṭ) — a crisis; a danger}

\label{lesson:MR-C295-sankat}



\begin{quote}
Before the new one: say the Marathi for a disaster, then the Marathi for a calamity.
\end{quote}

\subsection*{You'll want to know: \mr{संकट}}
\textbf{\mr{संकट}} — \emph{saṅkaṭ} — \textquotedblleft{}a crisis; a danger\textquotedblright{}.

\begin{quote}
\mr{गावावर} \mr{पाण्याचं} \mr{मोठं} \mr{संकट} \mr{आलं}.
\end{quote}

\begin{quote}
\emph{gāvāvar pāṇyācaṁ moṭhaṁ saṅkaṭ ālaṁ.}
\end{quote}

\begin{quote}
A great water crisis struck the village.
\end{quote}

Neuter.

\mr{संकट} \mr{येणे}, a crisis comes; \mr{संकटात}, in trouble.

\begin{cousinweb}
\textbf{\mr{संकट}} (\emph{saṅkaṭ}) is a tatsama, taken from Sanskrit \emph{saṅkaṭa}, 'narrow; a strait, a difficulty' (Monier-Williams).
\end{cousinweb}

\begin{grammarlens}[title={What you've built}]
One more word for telling what happened.
\end{grammarlens}

\subsection*{Your turn}
\begin{itemize}
\raggedright
  \item \emph{Say it:} \emph{saṅkaṭ}
  \item \emph{Say it:} \emph{saṅkaṭ}, once more
  \item \emph{Say it:} \emph{gāvāvar pāṇyācaṁ moṭhaṁ saṅkaṭ ālaṁ.} — A great water crisis struck the village.
  \item \emph{Recall:} say the Marathi for a disaster, then the Marathi for a calamity, then say \emph{saṅkaṭ} again
\end{itemize}

\begin{tcolorbox}[breakable,colback=teal!4,colframe=teal!35!black,title={Before you move on}]
What does \mr{संकट} mean? (\textquotedblleft{}A crisis; a danger\textquotedblright{}.) Say it once more.
\end{tcolorbox}
\section[\texorpdfstring{\mr{भीषण} (bhīṣaṇ)}{भीषण (bhīṣaṇ)}]{\mr{भीषण} (bhīṣaṇ) — dreadful; terrible}

\label{lesson:MR-C295-bhisan}



\begin{quote}
Before the new one: say the Marathi for a calamity, then the Marathi for a crisis.
\end{quote}

\subsection*{You'll want to know: \mr{भीषण}}
\textbf{\mr{भीषण}} — \emph{bhīṣaṇ} — \textquotedblleft{}dreadful; terrible\textquotedblright{}.

\begin{quote}
\mr{महामार्गावर} \mr{काल} \mr{भीषण} \mr{अपघात} \mr{झाला}.
\end{quote}

\begin{quote}
\emph{mahāmārgāvar kāl bhīṣaṇ apghāt jhālā.}
\end{quote}

\begin{quote}
There was a terrible accident on the highway yesterday.
\end{quote}

\mr{भीषण} keeps one form: \mr{भीषण} \mr{आग}, \mr{भीषण} \mr{अपघात}.

\begin{cousinweb}
\textbf{\mr{भीषण}} (\emph{bhīṣaṇ}) is a tatsama, taken from Sanskrit \emph{bhīṣaṇa}, 'terrifying' (Monier-Williams).
\end{cousinweb}

\begin{grammarlens}[title={What you've built}]
One more word for telling what happened.
\end{grammarlens}

\subsection*{Your turn}
\begin{itemize}
\raggedright
  \item \emph{Say it:} \emph{bhīṣaṇ}
  \item \emph{Say it:} \emph{bhīṣaṇ}, once more
  \item \emph{Say it:} \emph{mahāmārgāvar kāl bhīṣaṇ apghāt jhālā.} — There was a terrible accident on the highway yesterday.
  \item \emph{Recall:} say the Marathi for a calamity, then the Marathi for a crisis, then say \emph{bhīṣaṇ} again
\end{itemize}

\begin{tcolorbox}[breakable,colback=teal!4,colframe=teal!35!black,title={Before you move on}]
What does \mr{भीषण} mean? (\textquotedblleft{}Dreadful; terrible\textquotedblright{}.) Say it once more.
\end{tcolorbox}
\section[(dialogue)]{Looking back: incidents, witnesses and the police}

\label{lesson:MR-R295-review-288-295}



\begin{quote}
Say the Marathi for an event and an accident.
\end{quote}

\subsection*{Your turn}
Say these aloud:

\begin{itemize}
\raggedright
  \item an event, an accident, a theft, to sink, minor
  \item a witness, a crime, a robbery, a pickpocket, to cheat
  \item a constable, to arrest, to inquire, a prison, a prisoner
  \item to slip away, a rescue, safe and sound, missing, to comb through
  \item a spark, to flare up, a tongue of flame, ash, a fire brigade
  \item an earthquake, to shake, to collapse, to cave in, a landslide
  \item a flood, to collect, relief work, a volunteer, to get stuck
  \item damage, a disaster, a calamity, a crisis, dreadful
\end{itemize}

\begin{tcolorbox}[breakable,colback=teal!4,colframe=teal!35!black,title={Before you move on}]
An event? (\textbf{\mr{घटना}}, \emph{ghaṭnā}.) An accident? (\textbf{\mr{अपघात}}, \emph{apghāt}.) A theft? (\textbf{\mr{चोरी}}, \emph{corī}.) To sink? (\textbf{\mr{बुडणे}}, \emph{buḍṇe}.) Minor? (\textbf{\mr{किरकोळ}}, \emph{kirkoḷ}.) A witness? (\textbf{\mr{साक्षीदार}}, \emph{sākṣīdār}.) A crime? (\textbf{\mr{गुन्हा}}, \emph{gunhā}.) A robbery? (\textbf{\mr{दरोडा}}, \emph{daroḍā}.) A pickpocket? (\textbf{\mr{पाकीटमार}}, \emph{pākīṭmār}.) To cheat? (\textbf{\mr{फसवणे}}, \emph{phasavṇe}.) A constable? (\textbf{\mr{हवालदार}}, \emph{havāldār}.) To arrest? (\textbf{\mr{अटक} \mr{करणे}}, \emph{aṭak karṇe}.) To inquire? (\textbf{\mr{चौकशी} \mr{करणे}}, \emph{caukśī karṇe}.) A prison? (\textbf{\mr{तुरुंग}}, \emph{turuṅg}.) A prisoner? (\textbf{\mr{कैदी}}, \emph{kaidī}.) To slip away? (\textbf{\mr{निसटणे}}, \emph{nisaṭṇe}.) A rescue? (\textbf{\mr{सुटका}}, \emph{suṭkā}.) Safe and sound? (\textbf{\mr{सुखरूप}}, \emph{sukhrūp}.) Missing? (\textbf{\mr{बेपत्ता}}, \emph{bepattā}.) To comb through? (\textbf{\mr{धुंडाळणे}}, \emph{dhuṇḍāḷṇe}.) A spark? (\textbf{\mr{ठिणगी}}, \emph{ṭhiṇgī}.) To flare up? (\textbf{\mr{भडकणे}}, \emph{bhaḍakṇe}.) A tongue of flame? (\textbf{\mr{ज्वाळा}}, \emph{jvāḷā}.) Ash? (\textbf{\mr{राख}}, \emph{rākh}.) A fire brigade? (\textbf{\mr{अग्निशमन} \mr{दल}}, \emph{agniśaman dal}.) An earthquake? (\textbf{\mr{भूकंप}}, \emph{bhūkamp}.) To shake? (\textbf{\mr{हादरणे}}, \emph{hādarṇe}.) To collapse? (\textbf{\mr{कोसळणे}}, \emph{kosaḷṇe}.) To cave in? (\textbf{\mr{खचणे}}, \emph{khacṇe}.) A landslide? (\textbf{\mr{दरड}}, \emph{daraḍ}.) A flood? (\textbf{\mr{पूर}}, \emph{pūr}.) To collect? (\textbf{\mr{साचणे}}, \emph{sācṇe}.) Relief work? (\textbf{\mr{मदतकार्य}}, \emph{madatkārya}.) A volunteer? (\textbf{\mr{स्वयंसेवक}}, \emph{svayaṁsevak}.) To get stuck? (\textbf{\mr{अडकणे}}, \emph{aḍakṇe}.) Damage? (\textbf{\mr{नुकसान}}, \emph{nuksān}.) A disaster? (\textbf{\mr{दुर्घटना}}, \emph{durghaṭnā}.) A calamity? (\textbf{\mr{आपत्ती}}, \emph{āpattī}.) A crisis? (\textbf{\mr{संकट}}, \emph{saṅkaṭ}.) Dreadful? (\textbf{\mr{भीषण}}, \emph{bhīṣaṇ}.)
\end{tcolorbox}
